\documentclass[aspectratio=169, 10pt]{beamer}
% ==============================================================================
% Quantum Computing 12-Hour Course: Common LaTeX Beamer Preamble
% ==============================================================================
\usepackage{amsmath, amssymb, amsfonts, mathtools}
\usepackage{booktabs}
\usepackage{kotex}
\usepackage{tikz}
\usetikzlibrary{arrows.meta, positioning, shapes.geometric, calc}

% Syntax highlighting & listings
\usepackage{listings}
\usepackage{xcolor}

\definecolor{qcDarkBlue}{RGB}{20, 35, 60}
\definecolor{qcAccentBlue}{RGB}{0, 114, 206}
\definecolor{qcTeal}{RGB}{0, 168, 150}
\definecolor{qcLightGray}{RGB}{246, 248, 251}
\definecolor{qcCodeBg}{RGB}{242, 244, 247}
\definecolor{qcKeyword}{RGB}{140, 30, 120}
\definecolor{qcString}{RGB}{30, 120, 50}
\definecolor{qcComment}{RGB}{110, 110, 120}

\lstset{
    backgroundcolor=\color{qcCodeBg},
    basicstyle=\ttfamily\scriptsize,
    keywordstyle=\color{qcKeyword}\bfseries,
    stringstyle=\color{qcString},
    commentstyle=\color{qcComment},
    breaklines=true,
    showstringspaces=false,
    tabsize=4,
    frame=single,
    rulecolor=\color{qcAccentBlue!30},
    numbers=left,
    numberstyle=\tiny\color{qcComment},
    numbersep=6pt,
    xleftmargin=10pt,
    xrightmargin=5pt
}

% Beamer Theme Settings
\usefonttheme[onlymath]{serif}
\setbeamertemplate{navigation symbols}{}
\setbeamercolor{palette primary}{bg=qcDarkBlue, fg=white}
\setbeamercolor{palette secondary}{bg=qcAccentBlue, fg=white}
\setbeamercolor{palette tertiary}{bg=qcTeal, fg=white}
\setbeamercolor{structure}{fg=qcAccentBlue}
\setbeamercolor{titlelike}{parent=palette primary}
\setbeamercolor{frametitle}{bg=qcDarkBlue, fg=white}
\setbeamercolor{framesubtitle}{bg=qcDarkBlue, fg=white!85}
\setbeamerfont{frametitle}{size=\normalsize, series=\bfseries}
\setbeamerfont{framesubtitle}{size=\scriptsize, series=\normalfont}

\makeatletter
\setbeamertemplate{frametitle}{%
    \nointerlineskip%
    \begin{beamercolorbox}[wd=\paperwidth,sep=4pt,leftskip=8pt]{frametitle}%
        {\usebeamerfont{frametitle}\insertframetitle}%
        \ifx\insertframesubtitle\@empty%
        \else%
            \par\vspace{1pt}%
            {\usebeamerfont{framesubtitle}\usebeamercolor[fg]{framesubtitle}\insertframesubtitle}%
        \fi%
    \end{beamercolorbox}%
}
\makeatother

% Native Beamer Blocks (Rounded, theme-matched, 100% pause-compatible)
\setbeamertemplate{blocks}[rounded][shadow=false]

% 1. Standard block (formerly qcblock)
\setbeamercolor{block title}{fg=white, bg=qcAccentBlue}
\setbeamercolor{block body}{fg=qcDarkBlue, bg=qcLightGray}
\setbeamerfont{block title}{size=\footnotesize, series=\bfseries}

% 2. Alertblock (formerly qcalert)
\setbeamercolor{block title alerted}{fg=white, bg=red!75!black}
\setbeamercolor{block body alerted}{fg=black, bg=red!5}
\setbeamerfont{block title alerted}{size=\footnotesize, series=\bfseries}

% 3. Exampleblock (formerly qctip)
\setbeamercolor{block title example}{fg=white, bg=qcTeal}
\setbeamercolor{block body example}{fg=black, bg=green!5}
\setbeamerfont{block title example}{size=\footnotesize, series=\bfseries}

% Tight vertical block spacing
\addtobeamertemplate{block begin}{\vspace{-3pt}}{}
\addtobeamertemplate{block end}{}{\vspace{-4pt}}
\addtobeamertemplate{block alerted begin}{\vspace{-3pt}}{}
\addtobeamertemplate{block alerted end}{}{\vspace{-4pt}}
\addtobeamertemplate{block example begin}{\vspace{-3pt}}{}
\addtobeamertemplate{block example end}{}{\vspace{-4pt}}

% Backward compatibility aliases
\newenvironment{qcblock}[1]{\begin{block}{#1}}{\end{block}}
\newenvironment{qcalert}[1]{\begin{alertblock}{#1}}{\end{alertblock}}
\newenvironment{qctip}[1]{\begin{exampleblock}{#1}}{\end{exampleblock}}

% Quantum Dirac Notation Macros
\newcommand{\ket}[1]{|#1\rangle}
\newcommand{\bra}[1]{\langle#1|}
\newcommand{\braket}[2]{\langle#1|#2\rangle}
\newcommand{\ketbra}[2]{|#1\rangle\langle#2|}
\newcommand{\normsq}[1]{|#1|^2}

\title[양자컴퓨팅 부트캠프 2일차]{양자컴퓨팅 이론\&실습 4주 부트캠프}
\subtitle{2일차: 다중 큐비트와 양자 얽힘(Entanglement)}
\author{강사: 서상협 ((주)큐렌스)}
\date{2026. 11. 18 (수)}
\begin{document}

% Slide 1
\frame{\titlepage}

% Slide 2
\begin{frame}{2일차 핵심 질문: 독립된 계에서 상호작용 계로}
    \begin{center}
        \large \textbf{"큐비트가 결합하면 상태공간의 차원은 왜 지수적으로 팽창($2^n$)하며,\\독립된 상태로 분리할 수 없는 '얽힘'은 어떻게 만들어지는가?"}
    \end{center}
    \vspace{0.2cm}
    \begin{qcblock}{2일차 학습의 관점: 독립된 상태에서 '비분리 얽힘 상태'로}
        \begin{itemize}
            \item 고전 컴퓨터의 50비트는 독립된 2진 스위치 50개에 불과합니다.
            \item 큐비트가 결합하면 상태를 기술하는 힐베르트 공간의 차원이 \textbf{$2^n$으로 지수적 팽창}을 합니다.
            \item 두 큐비트가 상호작용하여 개별 상태의 텐서곱으로 분해할 수 없는 \textbf{'얽힘(Entanglement)'} 상태가 탄생합니다.
            \item 텐서곱($\otimes$)과 2큐비트 유니터리 연산(CNOT)의 선형대수학적 구조를 손계산과 NumPy로 확인합니다.
        \end{itemize}
    \end{qcblock}
\end{frame}

% Slide 3
\begin{frame}{2일차 세부 운영안}
    \begin{center}
    \resizebox{0.95\textwidth}{!}{%
    \begin{tabular}{clcp{7.5cm}}
    \toprule
    \textbf{구간} & \textbf{주제} & \textbf{시간} & \textbf{학습자 인지 궤적 및 활동} \\
    \midrule
    \textbf{1부} & \textbf{다중 큐비트 이론 및 손계산 집중} & 75분 & $2^n$ 차원 폭발 $\to$ 텐서곱($\otimes$) 대수 $\to$ CNOT 게이트 $\to$ \textbf{[손계산 5단계] Bell 회로 유도 및 얽힘 증명} \\
    \midrule
    \textbf{휴식} & \textbf{휴식 및 실습 환경 점검} & 20분 & 재충전, 1부 손계산 개별 질의응답, Python/NumPy 실습 환경 점검 \\
    \midrule
    \textbf{2부} & \textbf{NumPy 다중 큐비트 인터랙티브 실습} & 75분 & No-Signaling $\to$ \textbf{[Lab 1--3] 텐서곱 $\cdot$ CNOT $\cdot$ 얽힘 판별식} $\to$ \textbf{[2부 챌린지] 4대 Bell 상태 합성} $\to$ 부분 측정 붕괴 \\
    \bottomrule
    \end{tabular}%
    }
    \end{center}
\end{frame}

% Slide 4
\begin{frame}{1일차 복습: 단일 큐비트에서 다중 큐비트로의 확장}
    \begin{qcblock}{1일차에 배운 단일 큐비트 시스템 ($\mathbb{C}^2$)}
        \begin{itemize}
            \item 2차원 복소 힐베르트 공간: $\ket{\psi} = \alpha\ket{0} + \beta\ket{1}$ ($|\alpha|^2 + |\beta|^2 = 1$)
            \item 1큐비트 게이트: $2 \times 2$ 유니터리 행렬 ($X, Y, Z, H, S, T$)
            \item 상태의 기하학: 블로흐 구(Bloch Sphere) 위의 한 점 $(\theta, \phi)$
        \end{itemize}
    \end{qcblock}
    \vspace{0.2cm}
    \begin{qcalert}{오늘 마주할 2큐비트 이상의 복합 시스템 ($\mathbb{C}^2 \otimes \mathbb{C}^2 \cong \mathbb{C}^4$)}
        \begin{itemize}
            \item 큐비트가 2개로 늘어나면 기저 상태가 4개(`00`, `01`, `10`, `11`)로 늘어납니다.
            \item 더 이상 블로흐 구 2개만으로는 두 큐비트의 상호작용과 얽힘을 표현할 수 없습니다!
            \item 두 큐비트가 결합하는 수학적 연산자, \textbf{텐서곱(Tensor Product, $\otimes$)}이 필요한 이유입니다.
        \end{itemize}
    \end{qcalert}
\end{frame}

% Slide 5
\begin{frame}{상태공간 차원의 차이: 고전 비트 50개 vs 큐비트 50개}
    \begin{qcblock}{고전 비트 50개의 상태}
        \begin{itemize}
            \item 각 비트는 0 또는 1의 확정된 값 하나를 가집니다.
            \item 예: `11010010101...` $\implies$ 50개의 비트, 즉 \textbf{약 6.25 바이트}에 불과합니다.
            \item 최신 컴퓨터 메모리(수십 GB)와 비교해도 극소량의 데이터에 해당합니다.
        \end{itemize}
    \end{qcblock}
    \vspace{0.2cm}
    \begin{qcalert}{큐비트 50개의 상태}
        \begin{itemize}
            \item 큐비트 50개는 50개의 숫자로 표현되지 않습니다!
            \item 50개의 큐비트가 동시에 가질 수 있는 모든 2진수 경우의 수($00\dots0$부터 $11\dots1$까지)마다 \textbf{각각 고유한 복소수 진폭}이 붙습니다.
            \item 필요한 복소수 개수: \textbf{$2^{50} \approx 1,125,899,906,842,624$개 (약 1,125조 개!)}
        \end{itemize}
    \end{qcalert}
\end{frame}

% Slide 6
\begin{frame}{차원의 폭발: 왜 슈퍼컴퓨터의 RAM이 터지는가?}
    \begin{center}
    \resizebox{0.92\textwidth}{!}{%
    \begin{tabular}{ccl}
    \toprule
    \textbf{큐비트 수 ($n$)} & \textbf{필요한 진폭 개수 ($2^n$)} & \textbf{필요한 메모리 용량 (복소수 16바이트 기준)} \\
    \midrule
    1개 & $2^1 = 2$개 & 32 바이트 (스마트워치 수준) \\
    2개 & $2^2 = 4$개 & 64 바이트 \\
    10개 & $2^{10} = 1,024$개 & 16 킬로바이트 (KB) \\
    20개 & $2^{20} \approx 100\text{만}$개 & 16 메가바이트 (MB, 노트북 가능) \\
    30개 & $2^{30} \approx 10\text{억}$개 & 16 기가바이트 (GB, 고성능 워크스테이션 한계) \\
    \textbf{50개} & \textbf{$2^{50} \approx 1,125\text{조}$개} & \textbf{16 페타바이트 (PB, 전 세계 최고 슈퍼컴퓨터 메모리 초과!)} \\
    \bottomrule
    \end{tabular}%
    }
    \end{center}
    \vspace{0.2cm}
    \begin{qcblock}{결론}
        양자컴퓨터가 50큐비트만 도달해도 슈퍼컴퓨터가 흉내 낼 수 없는 이유는,\\
        \textbf{"이 지수적인 $2^n$차원 공간의 복소수들을 슈퍼컴퓨터 메모리에 담을 수조차 없기 때문"}입니다.
    \end{qcblock}
\end{frame}

% Slide 7
\begin{frame}{왜 양자 시뮬레이션은 고전 컴퓨터에서 불가능한가?}
    \begin{qcblock}{선형 대수학적 병목 (Linear Algebra Bottleneck)}
        \begin{itemize}
            \item 고전 알고리즘에서 변수가 1개 늘어나면 배열 크기는 $+1$ 됩니다 (선형 증가).
            \item 양자 시스템에서는 큐비트가 1개 늘어날 때마다 상태벡터 길이가 \textbf{$\times 2$ (2배)}가 됩니다.
            \item 50큐비트 게이트 연산은 \textbf{$2^{50} \times 2^{50}$ 크기의 복소 행렬}을 곱하는 것과 같습니다.
        \end{itemize}
    \end{qcblock}
    \vspace{0.2cm}
    \begin{qctip}{리처드 파인만(Richard Feynman)의 1981년 통찰}
        \begin{center}
            \textit{"Nature isn't classical, dammit, and if you want to make a simulation of nature,\\you'd better make it quantum mechanical!"}
        \end{center}
        자연계의 양자 현상을 컴퓨터로 모사하려면, 컴퓨터 자체가 양자역학적으로 동작해야 합니다.
    \end{qctip}
\end{frame}

% Slide 8
\begin{frame}{2큐비트 공간의 4대 기준점 (계산 기저)}
    큐비트가 2개가 되면, 고전 비트의 4가지 조합(`00`, `01`, `10`, `11`)이 기저 벡터가 됩니다:
    \vspace{0.2cm}
    \begin{qcblock}{4차원 열벡터 표현}
        $$\ket{00} = \begin{pmatrix} 1 \\ 0 \\ 0 \\ 0 \end{pmatrix}, \quad
          \ket{01} = \begin{pmatrix} 0 \\ 1 \\ 0 \\ 0 \end{pmatrix}, \quad
          \ket{10} = \begin{pmatrix} 0 \\ 0 \\ 1 \\ 0 \end{pmatrix}, \quad
          \ket{11} = \begin{pmatrix} 0 \\ 0 \\ 0 \\ 1 \end{pmatrix}$$
    \end{qcblock}
    \vspace{0.2cm}
    \begin{qctip}{임의의 2큐비트 중첩 상태}
        $$\ket{\psi} = c_{00}\ket{00} + c_{01}\ket{01} + c_{10}\ket{10} + c_{11}\ket{11} = \begin{pmatrix} c_{00} \\ c_{01} \\ c_{10} \\ c_{11} \end{pmatrix}$$
        전체 확률 합 1 조건: $|c_{00}|^2 + |c_{01}|^2 + |c_{10}|^2 + |c_{11}|^2 = 1$
    \end{qctip}
\end{frame}

% Slide 9
\begin{frame}{[기정사실화 안내] 텐서곱($\otimes$)의 대수적 정의}
    \begin{qcblock}{수학 기호 $\otimes$에 겁먹지 마세요!}
        \begin{itemize}
            \item 교재나 논문에 자주 등장하는 $\ket{0} \otimes \ket{0} = \ket{00}$ 기호는 \textbf{'텐서곱(Tensor Product)'}이라 부릅니다.
            \item 고등 수학에서는 이를 복잡한 텐서 대수로 정의하지만, 양자컴퓨팅에서는 아주 단순합니다:
            $$\begin{pmatrix} a \\ b \end{pmatrix} \otimes \begin{pmatrix} c \\ d \end{pmatrix} = \begin{pmatrix} a \cdot \begin{pmatrix} c \\ d \end{pmatrix} \\ b \cdot \begin{pmatrix} c \\ d \end{pmatrix} \end{pmatrix} = \begin{pmatrix} ac \\ ad \\ bc \\ bd \end{pmatrix}$$
            \item \textbf{규칙}: "앞 벡터의 각 성분에 뒤 벡터 전체를 곱해서 세로로 길게 이어 붙이는 블록 쌓기!"
        \end{itemize}
    \end{qcblock}
\end{frame}

% Slide 10
\begin{frame}{텐서곱 상세 전개 예제: 2차원 벡터의 결합}
    두 단일 큐비트 상태 $\ket{0} = \begin{pmatrix} 1 \\ 0 \end{pmatrix}$과 $\ket{1} = \begin{pmatrix} 0 \\ 1 \end{pmatrix}$의 텐서곱을 전개해 봅시다:
    \vspace{0.2cm}
    \begin{qcblock}{$\ket{01} = \ket{0} \otimes \ket{1}$의 단계별 계산}
        $$\ket{0} \otimes \ket{1} = \begin{pmatrix} 1 \\ 0 \end{pmatrix} \otimes \begin{pmatrix} 0 \\ 1 \end{pmatrix} = \begin{pmatrix} 1 \cdot \begin{pmatrix} 0 \\ 1 \end{pmatrix} \\ 0 \cdot \begin{pmatrix} 0 \\ 1 \end{pmatrix} \end{pmatrix} = \begin{pmatrix} 1 \cdot 0 \\ 1 \cdot 1 \\ 0 \cdot 0 \\ 0 \cdot 1 \end{pmatrix} = \begin{pmatrix} 0 \\ 1 \\ 0 \\ 0 \end{pmatrix} = \ket{01}$$
    \end{qcblock}
    \vspace{0.2cm}
    \begin{qctip}{$\ket{10} = \ket{1} \otimes \ket{0}$의 단계별 계산}
        $$\ket{1} \otimes \ket{0} = \begin{pmatrix} 0 \\ 1 \end{pmatrix} \otimes \begin{pmatrix} 1 \\ 0 \end{pmatrix} = \begin{pmatrix} 0 \cdot \begin{pmatrix} 1 \\ 0 \end{pmatrix} \\ 1 \cdot \begin{pmatrix} 1 \\ 0 \end{pmatrix} \end{pmatrix} = \begin{pmatrix} 0 \cdot 1 \\ 0 \cdot 0 \\ 1 \cdot 1 \\ 1 \cdot 0 \end{pmatrix} = \begin{pmatrix} 0 \\ 0 \\ 1 \\ 0 \end{pmatrix} = \ket{10}$$
        두 벡터의 순서가 바뀌면 결과 벡터의 1이 위치하는 인덱스도 완전히 달라집니다!
    \end{qctip}
\end{frame}

% Slide 11
\begin{frame}{행렬의 텐서곱: 2큐비트 유니터리 연산자 구성}
    단일 큐비트 게이트 2개를 각각의 큐비트에 동시에 적용할 때, 전체 연산자 행렬은 \textbf{게이트 행렬들의 텐서곱($4 \times 4$)}으로 만들어집니다:
    \vspace{0.2cm}
    \begin{qcblock}{$A \otimes B$의 블록 행렬 정의}
        $$A \otimes B = \begin{pmatrix} a_{11} & a_{12} \\ a_{21} & a_{22} \end{pmatrix} \otimes B = \begin{pmatrix} a_{11}B & a_{12}B \\ a_{21}B & a_{22}B \end{pmatrix}$$
    \end{qcblock}
    \vspace{0.2cm}
    \begin{qctip}{예시: 1번 큐비트에만 $H$를 적용하고 2번 큐비트는 가만히 둘 때 ($H \otimes I$)}
        $$H \otimes I = \frac{1}{\sqrt{2}}\begin{pmatrix} 1 & 1 \\ 1 & -1 \end{pmatrix} \otimes \begin{pmatrix} 1 & 0 \\ 0 & 1 \end{pmatrix} = \frac{1}{\sqrt{2}}\begin{pmatrix} 1 \cdot I & 1 \cdot I \\ 1 \cdot I & -1 \cdot I \end{pmatrix} = \frac{1}{\sqrt{2}}\begin{pmatrix} 1 & 0 & 1 & 0 \\ 0 & 1 & 0 & 1 \\ 1 & 0 & -1 & 0 \\ 0 & 1 & 0 & -1 \end{pmatrix}$$
    \end{qctip}
\end{frame}

% Slide 12
\begin{frame}{텐서곱의 핵심 성질: 혼합곱 공식(Mixed-Product Property)}
    양자 회로 계산을 비약적으로 단순화해 주는 가장 중요한 대수적 정리입니다:
    \vspace{0.2cm}
    \begin{qcalert}{혼합곱 공식 (Mixed-Product Property)}
        $$(A \otimes B)(C \otimes D) = (AC) \otimes (BD)$$
        "텐서곱을 한 뒤 행렬을 곱하는 것은, 각 게이트끼리 먼저 곱한 뒤 텐서곱을 취하는 것과 같다!"
    \end{qcalert}
    \vspace{0.2cm}
    \begin{qcblock}{행렬과 벡터의 결합 형태}
        $$(A \otimes B)(\ket{\psi_1} \otimes \ket{\psi_2}) = (A\ket{\psi_1}) \otimes (B\ket{\psi_2})$$
        \begin{itemize}
            \item $4 \times 4$ 행렬 전체를 거대하게 계산하지 않고도,
            \item 1번 큐비트 연산 $A\ket{\psi_1}$과 2번 큐비트 연산 $B\ket{\psi_2}$를 각각 독립적으로 계산한 뒤 텐서곱할 수 있습니다!
        \end{itemize}
    \end{qcblock}
\end{frame}

% Slide 13
\begin{frame}{서로 대화하는 2큐비트 게이트: CNOT (제어-NOT)}
    1큐비트 게이트($H, X, Z$)는 각 큐비트가 혼자 도는 연산이었습니다.
    양자컴퓨터가 강력해지려면 \textbf{두 큐비트가 서로 영향을 주고받는 제어 연산}이 필수적입니다.
    \vspace{0.2cm}
    \begin{qcalert}{CNOT (Controlled-NOT) 게이트의 동작 규칙}
        \begin{itemize}
            \item 1번 큐비트(\textbf{제어 큐비트}): 상태를 바꾸지 않고 신호만 줍니다.
            \item 2번 큐비트(\textbf{타깃 큐비트}): 제어 큐비트의 상태에 따라 동작합니다.
        \end{itemize}
        \begin{center}
            \LARGE \color{qcAccentBlue} \textbf{"제어 큐비트가 1이면, 타깃 큐비트를 NOT(뒤집기)한다!"}
        \end{center}
    \end{qcalert}
\end{frame}

% Slide 14
\begin{frame}{CNOT 게이트의 진리표와 회로 기호}
    \begin{center}
    \begin{tikzpicture}[thick, scale=1.0]
        \draw (0, 1.0) node[left] {$q_0$ (제어)} -- (4.0, 1.0) node[right] {$q_0$};
        \draw (0, 0.0) node[left] {$q_1$ (타깃)} -- (4.0, 0.0) node[right] {$q_1 \oplus q_0$};
        \node[circle, fill=black, inner sep=2.5pt] (ctrl) at (2.0, 1.0) {};
        \node[circle, draw, thick, fill=white, minimum size=0.55cm, inner sep=0pt] (targ) at (2.0, 0.0) {};
        \draw[thick] (targ.north) -- (targ.south);
        \draw[thick] (targ.west) -- (targ.east);
        \draw (ctrl) -- (targ.north);
    \end{tikzpicture}
    \end{center}
    \vspace{0.1cm}
    \begin{center}
    \begin{tabular}{ccc}
    \toprule
    \textbf{입력 상태 ($\ket{q_0 q_1}$)} & \textbf{제어 비트 ($q_0$)} & \textbf{출력 상태} \\
    \midrule
    $\ket{00}$ & $0$ (뒤집지 않음) & $\ket{00}$ \\
    $\ket{01}$ & $0$ (뒤집지 않음) & $\ket{01}$ \\
    $\ket{10}$ & $1$ (뒤집음: $0 \to 1$) & $\ket{11}$ \\
    $\ket{11}$ & $1$ (뒤집음: $1 \to 0$) & $\ket{10}$ \\
    \bottomrule
    \end{tabular}
    \end{center}
\end{frame}

% Slide 15
\begin{frame}{CNOT의 $4 \times 4$ 행렬 구조 유도}
    4가지 계산 기저에 작용한 출력을 열벡터로 나열하면 CNOT 행렬이 완성됩니다:
    $$CX\ket{00} = \ket{00}, \quad CX\ket{01} = \ket{01}, \quad CX\ket{10} = \ket{11}, \quad CX\ket{11} = \ket{10}$$
    \vspace{0.2cm}
    \begin{qcblock}{CNOT의 $4 \times 4$ 유니터리 행렬 표현}
        $$CX = \begin{pmatrix}
        1 & 0 & 0 & 0 \\
        0 & 1 & 0 & 0 \\
        0 & 0 & 0 & 1 \\
        0 & 0 & 1 & 0
        \end{pmatrix} = \begin{pmatrix} I_2 & 0 \\ 0 & X \end{pmatrix}$$
        \begin{itemize}
            \item $q_0=0$일 때(위쪽 $2 \times 2$ 블록): 항등행렬 $I_2$가 작용하여 상태를 유지합니다.
            \item $q_0=1$일 때(아래쪽 $2 \times 2$ 블록): Pauli-$X$ 행렬이 작용하여 비트를 반전시킵니다.
        \end{itemize}
    \end{qcblock}
\end{frame}

% Slide 16
\begin{frame}{양자 얽힘(Entanglement) 만들기: Bell 회로 개요}
    이제 $H$ 게이트와 CNOT 게이트를 결합하여, 양자 알고리즘 계산 이점의 핵심 기반인 \textbf{얽힘(Bell 상태)}을 직접 만듭니다:
    \vspace{0.3cm}
    \begin{center}
    \begin{tikzpicture}[thick, scale=1.0]
        \draw (0, 1.0) node[left] {$q_0 : \ket{0}$} -- (6.0, 1.0) node[right] {?};
        \draw (0, 0.0) node[left] {$q_1 : \ket{0}$} -- (6.0, 0.0) node[right] {?};
        \node[draw, fill=qcAccentBlue!20, minimum size=0.7cm] at (1.6, 1.0) {$H$};
        \node[circle, fill=black, inner sep=2.5pt] (ctrl) at (4.0, 1.0) {};
        \node[circle, draw, thick, fill=white, minimum size=0.55cm, inner sep=0pt] (targ) at (4.0, 0.0) {};
        \draw[thick] (targ.north) -- (targ.south);
        \draw[thick] (targ.west) -- (targ.east);
        \draw (ctrl) -- (targ.north);
        \node at (0.7, -0.6) {\small [초기 상태]};
        \node at (2.8, -0.6) {\small [1단계: 중첩]};
        \node at (5.0, -0.6) {\small [2단계: CNOT]};
    \end{tikzpicture}
    \end{center}
    \vspace{0.2cm}
    \begin{qcblock}{워크시트 3 목표}
        초기 상태 $\ket{00}$에서 출발하여, 중간 상태 $\ket{\psi_1}$과 최종 얽힘 상태 $\ket{\Phi^+}$를 단계별로 직접 손계산합니다.
    \end{qcblock}
\end{frame}

% Slide 17
\begin{frame}{[손계산 1단계] 초기 상태와 1번 큐비트 Hadamard 작용}
    \begin{qcblock}{초기 상태의 텐서곱 표현}
        $$\ket{\psi_0} = \ket{0} \otimes \ket{0} = \begin{pmatrix} 1 \\ 0 \end{pmatrix} \otimes \begin{pmatrix} 1 \\ 0 \end{pmatrix} = \begin{pmatrix} 1 \\ 0 \\ 0 \\ 0 \end{pmatrix} = \ket{00}$$
    \end{qcblock}
    \vspace{0.2cm}
    \begin{qcalert}{1단계: $q_0$에 Hadamard ($H$) 게이트 적용}
        $$\ket{\psi_1} = (H \otimes I)\ket{\psi_0} = (H\ket{0}) \otimes (I\ket{0})$$
        $$= \left( \frac{\ket{0} + \ket{1}}{\sqrt{2}} \right) \otimes \ket{0} = \frac{1}{\sqrt{2}}(\ket{0} \otimes \ket{0}) + \frac{1}{\sqrt{2}}(\ket{1} \otimes \ket{0})$$
        $$\implies \ket{\psi_1} = \frac{\ket{00} + \ket{10}}{\sqrt{2}} = \frac{1}{\sqrt{2}}\begin{pmatrix} 1 \\ 0 \\ 1 \\ 0 \end{pmatrix}$$
    \end{qcalert}
\end{frame}

% Slide 18
\begin{frame}{[손계산 2단계] CNOT 게이트 통과와 조건부 반전}
    \begin{qcblock}{중간 상태 $\ket{\psi_1}$에 CNOT 게이트 작용}
        $$\ket{\psi_2} = CX \ket{\psi_1} = CX \left( \frac{\ket{00} + \ket{10}}{\sqrt{2}} \right) = \frac{1}{\sqrt{2}} CX\ket{00} + \frac{1}{\sqrt{2}} CX\ket{10}$$
    \end{qcblock}
    \vspace{0.2cm}
    \begin{qctip}{선형성에 따른 각 항별 CNOT 연산 결과}
        \begin{itemize}
            \item 첫 번째 항 $CX\ket{00}$: 제어 큐비트가 $0$이므로 $\implies \ket{00}$ 유지!
            \item 두 번째 항 $CX\ket{10}$: 제어 큐비트가 $1$이므로 타깃 비트가 $0 \to 1$ 반전 $\implies \ket{11}$로 변환!
        \end{itemize}
    \end{qctip}
    \vspace{0.2cm}
    $$\implies \ket{\psi_2} = \frac{\ket{00} + \ket{11}}{\sqrt{2}} = \ket{\Phi^+}$$
\end{frame}

% Slide 19
\begin{frame}{[손계산 3단계] $4 \times 4$ 행렬곱으로 수치 재검증}
    행렬과 열벡터의 직접 곱셈으로 동일한 결과가 도출되는지 확인합니다:
    \vspace{0.2cm}
    \begin{qcblock}{행렬-벡터 곱셈 전개}
        $$\ket{\psi_2} = \begin{pmatrix}
        1 & 0 & 0 & 0 \\
        0 & 1 & 0 & 0 \\
        0 & 0 & 0 & 1 \\
        0 & 0 & 1 & 0
        \end{pmatrix}
        \begin{pmatrix}
        1/\sqrt{2} \\ 0 \\ 1/\sqrt{2} \\ 0
        \end{pmatrix}
        = \begin{pmatrix}
        1 \cdot (1/\sqrt{2}) + 0 + 0 + 0 \\
        0 + 1 \cdot 0 + 0 + 0 \\
        0 + 0 + 0 \cdot (1/\sqrt{2}) + 1 \cdot 0 \\
        0 + 0 + 1 \cdot (1/\sqrt{2}) + 0
        \end{pmatrix}
        = \begin{pmatrix}
        1/\sqrt{2} \\ 0 \\ 0 \\ 1/\sqrt{2}
        \end{pmatrix}$$
    \end{qcblock}
    \vspace{0.2cm}
    \begin{center}
        \large \color{qcAccentBlue} \textbf{Dirac 표기법 전개와 $4 \times 4$ 행렬 연산이 완벽히 일치합니다!}
    \end{center}
\end{frame}

% Slide 20
\begin{frame}{[손계산 검증] Bell 상태의 측정 확률 분석}
    $$\ket{\Phi^+} = \frac{1}{\sqrt{2}}\ket{00} + 0\ket{01} + 0\ket{10} + \frac{1}{\sqrt{2}}\ket{11}$$
    \vspace{0.2cm}
    \begin{qcblock}{Born 규칙에 따른 4가지 측정 확률 계산}
        \begin{itemize}
            \item $P(00) = \left|\frac{1}{\sqrt{2}}\right|^2 = \frac{1}{2} = 50\%$
            \item $P(01) = |0|^2 = 0\%$ (절대 관측되지 않음)
            \item $P(10) = |0|^2 = 0\%$ (절대 관측되지 않음)
            \item $P(11) = \left|\frac{1}{\sqrt{2}}\right|^2 = \frac{1}{2} = 50\%$
        \end{itemize}
    \end{qcblock}
    \vspace{0.2cm}
    \begin{qcalert}{놀라운 물리적 결과}
        두 큐비트를 동시에 측정하면, 서로 다른 값(`01`, `10`)은 전혀 나오지 않고\\
        \textbf{오직 같은 값(`00` 또는 `11`)만 각각 $50\%$ 확률로 완벽하게 일치하여 관측}됩니다!
    \end{qcalert}
\end{frame}

% Slide 21
\begin{frame}{[손계산 증명] 왜 얽힘 상태는 분리될 수 없는가? (비분리성 증명)}
    \begin{qcblock}{귀류법 증명: $\ket{\Phi^+}$가 독립된 두 상태의 텐서곱이라 가정}
        $$\ket{\Phi^+} = (a\ket{0} + b\ket{1}) \otimes (c\ket{0} + d\ket{1}) = ac\ket{00} + ad\ket{01} + bc\ket{10} + bd\ket{11}$$
    \end{qcblock}
    \vspace{0.2cm}
    \begin{qcalert}{계수 비교를 통한 모순 도출}
        \begin{enumerate}
            \item $ac = 1/\sqrt{2} \implies a \neq 0$ 이고 $c \neq 0$
            \item $ad = 0 \implies a \neq 0$ 이므로 반드시 \textbf{$d = 0$}
            \item $bc = 0 \implies c \neq 0$ 이므로 반드시 \textbf{$b = 0$}
            \item 그런데 네 번째 항 $bd = 1/\sqrt{2}$이어야 하는데, $b=0, d=0$이므로 $bd = 0 \neq 1/\sqrt{2}$! \textbf{(모순!)}
        \end{enumerate}
    \end{qcalert}
    \vspace{0.1cm}
    \begin{center}
        \textbf{결론: Bell 상태 $\ket{\Phi^+}$는 단일 큐비트 상태들의 곱으로 절대 쪼갤 수 없는 '얽힘(Entangled)' 상태입니다!}
    \end{center}
\end{frame}

% Slide 22
\begin{frame}{1부 마무리 점검 및 중간 휴식 안내 (20분)}
    \begin{qcblock}{1부 핵심 달성도 체크}
        \begin{itemize}
            \item $\checkmark$ 고전 50비트와 큐비트 50개의 차이 ($2^n$ 차원 힐베르트 공간의 지수적 폭발)
            \item $\checkmark$ 텐서곱($\otimes$)과 혼합곱 공식 $(A \otimes B)(C \otimes D) = (AC) \otimes (BD)$
            \item $\checkmark$ 제어-NOT(CNOT) 게이트의 조건부 반전 규칙과 $4 \times 4$ 행렬 작용
            \item $\checkmark$ [손계산 워크시트] Bell 상태 $\frac{\ket{00}+\ket{11}}{\sqrt{2}}$ 유도 및 비분리성 증명 완료
        \end{itemize}
    \end{qcblock}
    \vspace{0.3cm}
    \begin{qctip}{20분간 휴식 및 2부 NumPy 실습 준비}
        \begin{itemize}
            \item 휴식 후 2부에서는 \textbf{NumPy를 활용한 다중 큐비트 행렬 연산 및 얽힘 시뮬레이션(75분)}이 진행됩니다.
            \item Jupyter Notebook에서 실습 파일(\texttt{day2\_numpy\_entanglement.ipynb})을 열어주세요.
        \end{itemize}
    \end{qctip}
\end{frame}

% Slide 23
\begin{frame}{2부 실습 개요: 선형대수학을 코드로 정복하기}
    \begin{center}
        \large \textbf{"손으로 유도한 수학을 파이썬 배열과 행렬곱으로 100\% 재현한다"}
    \end{center}
    \vspace{0.2cm}
    \begin{qcblock}{2부 인터랙티브 실습 로드맵 (75분)}
        \begin{enumerate}
            \item \textbf{[컴퓨팅 원리]} 얽힘과 양자 우위: 얽힘이 없으면 50큐비트도 스마트폰 1초 컷!
            \item \textbf{[알고리즘 연계]} 왜 Bell 상태인가? 다중 큐비트 알고리즘의 최소 연산 단위
            \item \textbf{[Lab 1]} `np.kron`을 이용한 텐서곱과 혼합곱 성질 수치 검증
            \item \textbf{[Lab 2]} $4 \times 4$ CNOT 행렬곱 연산과 손계산 결과 1:1 대조 (`assert`)
            \item \textbf{[Lab 3]} 선형대수학적 얽힘 판별식 ($c_{00}c_{11} - c_{01}c_{10} \neq 0$) 함수 구현
            \item \textbf{[2부 챌린지]} 4대 Bell 상태 ($\ket{\Phi^\pm}, \ket{\Psi^\pm}$) 전수 합성 및 직교성 검증
            \item \textbf{[Lab 4]} QPU 내부의 부분 측정(Partial Measurement)과 조건부 상태 제어
        \end{enumerate}
    \end{qcblock}
\end{frame}

% Slide 24
\begin{frame}{얽힘의 컴퓨터공학적 본질: 양자 우위의 유일한 원천}
    \begin{qcalert}{만약 50큐비트 회로에 '얽힘'이 발생하지 않는다면?}
        모든 큐비트가 독립된 곱상태로 유지된다면:
        $$\ket{\psi} = (\alpha_1\ket{0}+\beta_1\ket{1}) \otimes (\alpha_2\ket{0}+\beta_2\ket{1}) \otimes \dots \otimes (\alpha_{50}\ket{0}+\beta_{50}\ket{1})$$
        \begin{itemize}
            \item 50개 큐비트 상태를 저장하는 데 필요한 복소수는 $2^{50}$(1,125조 개)가 아닙니다.
            \item \textbf{오직 $2 \times 50 = 100$개의 복소수}만 있으면 메모리에 전부 담깁니다!
            \item 즉, 스마트폰 메모리 1KB로도 0.001초 만에 완벽하게 시뮬레이션됩니다.
        \end{itemize}
    \end{qcalert}
    \vspace{0.2cm}
    \begin{qcblock}{컴퓨터 관점에서의 결론}
        양자컴퓨터가 고전 슈퍼컴퓨터를 압도할 수 있는 유일한 원천은\\
        \textbf{"개별 상태로 쪼갤 수 없는 비분리 얽힘 상태(Entangled State)의 존재"} 때문입니다!
    \end{qcblock}
\end{frame}

% Slide 25
\begin{frame}{다중 큐비트 알고리즘의 빌딩 블록: 왜 Bell 회로인가?}
    \begin{qcblock}{왜 우리는 Bell 회로를 가장 먼저 배우는가?}
        \begin{itemize}
            \item 1큐비트 게이트($H, X, Z$)만으로는 큐비트 간에 데이터를 주고받을 수 없습니다.
            \item 큐비트 간 상호작용(CNOT, CZ)이 결합하는 순간, \textbf{모든 양자 회로는 얽힘을 형성}합니다.
            \item 3일차에 배울 **Grover 알고리즘의 오라클과 확산 연산자** 역시 큐비트들을 정밀하게 얽어 파동을 간섭시킴으로써 정답 확률을 극대화합니다.
        \end{itemize}
    \end{qcblock}
    \vspace{0.2cm}
    \begin{qctip}{프로그래밍의 Hello World}
        프로그래밍 언어를 배울 때 가장 먼저 `printf("Hello World");`를 출력하듯,\\
        \textbf{Bell 회로는 모든 다중 큐비트 양자 알고리즘을 구성하는 가장 기초적인 연산 단위}입니다.
    \end{qctip}
\end{frame}

% Slide 26
\begin{frame}[fragile]{[Lab 1-1] NumPy로 4차원 기저 벡터 정의하기}
    \begin{lstlisting}[language=Python]
import numpy as np

# 1. 단일 큐비트 기저 정의
ket0 = np.array([[1.0], [0.0]], dtype=complex)
ket1 = np.array([[0.0], [1.0]], dtype=complex)

# 2. np.kron을 이용한 4차원 계산 기저 생성
ket00 = np.kron(ket0, ket0)  # |00> = [1, 0, 0, 0]^T
ket01 = np.kron(ket0, ket1)  # |01> = [0, 1, 0, 0]^T
ket10 = np.kron(ket1, ket0)  # |10> = [0, 0, 1, 0]^T
ket11 = np.kron(ket1, ket1)  # |11> = [0, 0, 0, 1]^T

print("|00> 기저 벡터:\n", ket00.flatten().real)
print("|11> 기저 벡터:\n", ket11.flatten().real)
assert np.allclose(ket00[0], 1.0) and np.allclose(ket11[3], 1.0)
\end{lstlisting}
\end{frame}

% Slide 27
\begin{frame}[fragile]{[Lab 1-2] 게이트 텐서곱과 혼합곱 공식 검증}
    \begin{lstlisting}[language=Python]
# 단일 큐비트 유니터리 게이트
H = (1 / np.sqrt(2)) * np.array([[1, 1], [1, -1]], dtype=complex)
X = np.array([[0, 1], [1, 0]], dtype=complex)
Z = np.array([[1, 0], [0, -1]], dtype=complex)

# 혼합곱 공식: (A (x) B)(C (x) D) == (AC) (x) (BD)
# 좌변 (LHS): 4x4 행렬 2개를 각각 만든 후 행렬곱
lhs = np.kron(H, X) @ np.kron(Z, H)

# 우변 (RHS): 2x2 행렬곱을 먼저 수행한 후 텐서곱
rhs = np.kron(H @ Z, X @ H)

# 수치 검증
assert np.allclose(lhs, rhs), "혼합곱 공식이 불일치합니다!"
print("[Lab 1 성공] (H (x) X)(Z (x) H) == (HZ) (x) (XH) 일치 확인 완료!")
\end{lstlisting}
\end{frame}

% Slide 28
\begin{frame}[fragile]{[Lab 2-1] $4 \times 4$ CNOT 행렬과 $H \otimes I$ 구성}
    \begin{lstlisting}[language=Python]
# 1. 4x4 CNOT 행렬 정의
CNOT = np.array([
    [1, 0, 0, 0],
    [0, 1, 0, 0],
    [0, 0, 0, 1],
    [0, 0, 1, 0]
], dtype=complex)

# 2. q0에만 H를 적용하는 4x4 연산자 H (x) I
I2 = np.eye(2, dtype=complex)
H_I = np.kron(H, I2)

# 유니터리성 전수 검증: U^dag @ U == I_4
assert np.allclose(CNOT.conj().T @ CNOT, np.eye(4, dtype=complex))
assert np.allclose(H_I.conj().T @ H_I, np.eye(4, dtype=complex))
print("CNOT 및 H (x) I 연산자 유니터리성 검증 통과!")
\end{lstlisting}
\end{frame}

% Slide 29
\begin{frame}[fragile]{[Lab 2-2] 행렬곱(@)으로 Bell 상태 생성 및 손계산 대조}
    \begin{lstlisting}[language=Python]
# 1. 초기 상태 |00>에 H (x) I 작용
psi_step1 = H_I @ ket00      # (|00> + |10>) / sqrt(2)

# 2. CNOT 게이트 통과 -> Bell 상태 |Phi+>
psi_bell = CNOT @ psi_step1  # (|00> + |11>) / sqrt(2)

# 3. 1부 손계산 결과와 오차 1e-15 이내 1:1 대조
expected_bell = (1.0 / np.sqrt(2.0)) * (ket00 + ket11)
assert np.allclose(psi_bell, expected_bell), "손계산 결과와 수치가 일치하지 않습니다!"

print("[Lab 2 성공] 1부 손계산 Bell 상태와 100% 일치!")
print("계산된 상태벡터:\n", psi_bell.flatten().real)
\end{lstlisting}
\end{frame}

% Slide 30
\begin{frame}{선형대수학적 얽힘 판별식 (비분리성 판별 조건)}
    임의의 2큐비트 상태 $\ket{\psi} = c_{00}\ket{00} + c_{01}\ket{01} + c_{10}\ket{10} + c_{11}\ket{11}$의 계수 행렬 $M = \begin{pmatrix} c_{00} & c_{01} \\ c_{10} & c_{11} \end{pmatrix}$:
    \vspace{0.1cm}
    \begin{qcblock}{행렬식(Determinant)과 분리 가능성}
        $$\det(M) = c_{00}c_{11} - c_{01}c_{10}$$
        \begin{itemize}
            \item $\det(M) = 0$: \textbf{곱상태} $\implies$ $(\alpha\ket{0}+\beta\ket{1})\otimes(\gamma\ket{0}+\delta\ket{1})$로 인수분해 가능
            \item $\det(M) \neq 0$: \textbf{얽힘 상태} $\implies$ 독립 상태의 곱으로 인수분해 절대 불가
        \end{itemize}
    \end{qcblock}
    \vspace{0.1cm}
    \begin{qctip}{Bell 상태의 행렬식 값}
        $\ket{\Phi^+} \implies c_{00} = 1/\sqrt{2}, c_{01}=0, c_{10}=0, c_{11}=1/\sqrt{2}$
        $$\det = (1/\sqrt{2})(1/\sqrt{2}) - 0 \cdot 0 = 1/2 \neq 0 \implies \text{최대 얽힘(Maximal Entanglement)!}$$
    \end{qctip}
\end{frame}

% Slide 31
\begin{frame}[fragile]{[Lab 3] 얽힘 판별식 함수 구현 및 수치 검증}
    \begin{lstlisting}[language=Python]
def is_entangled(state_4d, tol=1e-9):
    c = state_4d.flatten()
    det = c[0] * c[3] - c[1] * c[2]
    entangled = not np.isclose(abs(det), 0.0, atol=tol)
    return entangled, float(abs(det))

# 1. Step 1 상태 |+0> 판별: 곱상태이므로 det == 0이어야 함
ent1, det1 = is_entangled(psi_step1)
assert not ent1 and det1 == 0.0, "Step 1 상태는 분리 가능한 곱상태입니다!"

# 2. 최종 Bell 상태 |Phi+> 판별: 얽힘 상태이므로 det == 0.5이어야 함
ent2, det2 = is_entangled(psi_bell)
assert ent2 and np.isclose(det2, 0.5), "Bell 상태는 최대 얽힘 상태입니다!"

print(f"[Lab 3 성공] 곱상태 det={det1:.1f} (분리 가능) / Bell 상태 det={det2:.3f} (최대 얽힘)")
\end{lstlisting}
\end{frame}

% Slide 32
\begin{frame}{4대 Bell 상태 (Bell Basis)의 전체 체계}
    초기 입력 4가지 기저에 동일한 회로 $CNOT \cdot (H \otimes I)$를 통과시키면 4대 Bell 기저가 완성됩니다:
    \vspace{0.2cm}
    \begin{center}
    \begin{tabular}{cclc}
    \toprule
    \textbf{입력 기저} & \textbf{중간 상태 ($\ket{\psi_1}$)} & \textbf{최종 Bell 상태 ($\ket{\psi_2}$)} & \textbf{기호} \\
    \midrule
    $\ket{00}$ & $\frac{\ket{00} + \ket{10}}{\sqrt{2}}$ & $\frac{\ket{00} + \ket{11}}{\sqrt{2}}$ & $\ket{\Phi^+}$ \\
    $\ket{01}$ & $\frac{\ket{01} + \ket{11}}{\sqrt{2}}$ & $\frac{\ket{01} + \ket{10}}{\sqrt{2}}$ & $\ket{\Psi^+}$ \\
    $\ket{10}$ & $\frac{\ket{00} - \ket{10}}{\sqrt{2}}$ & $\frac{\ket{00} - \ket{11}}{\sqrt{2}}$ & $\ket{\Phi^-}$ \\
    $\ket{11}$ & $\frac{\ket{01} - \ket{11}}{\sqrt{2}}$ & $\frac{\ket{01} - \ket{10}}{\sqrt{2}}$ & $\ket{\Psi^-}$ \\
    \bottomrule
    \end{tabular}
    \end{center}
    \vspace{0.2cm}
    \begin{qcblock}{4대 Bell 기저의 직교성 (Orthogonality)}
        이 4개의 얽힘 상태들은 4차원 복소 공간을 이루는 \textbf{새로운 정규직교 기저(Orthonormal Basis)}가 됩니다.
    \end{qcblock}
\end{frame}

% Slide 33
\begin{frame}[fragile]{[2부 챌린지] NumPy로 4대 Bell 상태 전수 합성 및 직교성 증명}
    \begin{lstlisting}[language=Python]
inputs = {"|00>": ket00, "|01>": ket01, "|10>": ket10, "|11>": ket11}
bell = {name: CNOT @ (H_I @ vec) for name, vec in inputs.items()}

# 4대 Bell 상태: |Phi+>, |Psi+>, |Phi->, |Psi->
phi_p, psi_p = bell["|00>"], bell["|01>"]
phi_m, psi_m = bell["|10>"], bell["|11>"]

# 4개 상태의 상호 직교성 전수 검증: <u|v> == 0
assert np.isclose(np.vdot(phi_p, phi_m), 0.0)
assert np.isclose(np.vdot(phi_p, psi_p), 0.0)
assert np.isclose(np.vdot(phi_p, psi_m), 0.0)
assert np.isclose(np.vdot(psi_p, phi_m), 0.0)
assert np.isclose(np.vdot(psi_p, psi_m), 0.0)
assert np.isclose(np.vdot(phi_m, psi_m), 0.0)
print("[2부 챌린지 성공] 4대 Bell 상태가 상호 직교 기저임을 전수 증명 완료!")
\end{lstlisting}
\end{frame}

% Slide 34
\begin{frame}{QPU 내부의 부분 측정: 양자 오류 정정(QEC)과 조건부 제어}
    원격 통신이 아니라, \textbf{실제 양자 프로세서(QPU) 칩 내부}에서 부분 측정이 왜 필수적일까요?
    \vspace{0.2cm}
    \begin{qcalert}{양자 계산의 딜레마와 해결책}
        \begin{itemize}
            \item 계산 도중 메인 큐비트를 직접 측정하면 계산 중이던 중첩과 얽힘이 파괴됩니다.
            \item \textbf{해결책}: 연산 큐비트 옆에 보조 큐비트(Ancilla)를 얽어놓고, \textbf{보조 큐비트만 부분 측정}합니다!
            \item 메인 데이터의 중첩은 유지한 채 오류 발생 여부(신드롬)만 감지하는 양자 오류 정정(QEC)의 기본 원리입니다.
        \end{itemize}
    \end{qcalert}
    \vspace{0.1cm}
    \begin{qcblock}{수학적 사영 연산자(Projector)}
        $P_0 = \ketbra{0}{0} \otimes I, \quad P_1 = \ketbra{1}{1} \otimes I \implies \ket{\psi_{q1|k}} = \frac{P_k \ket{\Phi^+}}{\sqrt{p_k}}$
    \end{qcblock}
\end{frame}

% Slide 35
\begin{frame}[fragile]{[Lab 4] 부분 측정 시뮬레이션 및 조건부 붕괴 함수 구현}
    \begin{lstlisting}[language=Python]
def measure_q0(state_4d):
    v = state_4d.flatten()
    p0 = float(abs(v[0])**2 + abs(v[1])**2)  # q0=0일 확률 (|00>, |01>)
    p1 = float(abs(v[2])**2 + abs(v[3])**2)  # q0=1일 확률 (|10>, |11>)
    state_q1_0 = np.array([v[0], v[1]], dtype=complex) / np.sqrt(p0)
    state_q1_1 = np.array([v[2], v[3]], dtype=complex) / np.sqrt(p1)
    return (p0, state_q1_0), (p1, state_q1_1)

(p0, s0), (p1, s1) = measure_q0(psi_bell)
print(f"q0=0 관측(확률 {p0*100:.0f}%): q1 상태는 즉시 |0> ->", s0.real)
print(f"q0=1 관측(확률 {p1*100:.0f}%): q1 상태는 즉시 |1> ->", s1.real)
assert np.allclose(s0, [1, 0]) and np.allclose(s1, [0, 1])
print("[Lab 4 성공] QPU 내부의 조건부 상태 붕괴 메커니즘 검증 완료!")
\end{lstlisting}
\end{frame}

% Slide 36
\begin{frame}{2일차 총정리: 다중 큐비트와 얽힘의 핵심 원리}
    \begin{qcblock}{오늘 정리한 3대 컴퓨터공학적 원리}
        \begin{enumerate}
            \item \textbf{상태공간의 지수적 확장 ($2^n$)}:
            고전 비트 $n$개의 상태는 $n$개지만, 큐비트 $n$개의 복합 상태는 $2^n$차원 복소 힐베르트 공간의 단위벡터로 기술됩니다.
            \item \textbf{2큐비트 유니터리 연산: CNOT}:
            제어 큐비트 상태에 따라 타깃 큐비트에 $X$를 조건부 적용하여 큐비트 간 상호작용을 형성합니다.
            \item \textbf{양자 얽힘: 양자 우위의 유일한 계산 자원}:
            얽힘이 없다면 50큐비트도 스마트폰 1KB로 시뮬레이션됩니다. 비분리 얽힘 상태($c_{00}c_{11} - c_{01}c_{10} \neq 0$)가 존재해야 비로소 고전 슈퍼컴퓨터를 압도하는 진정한 양자 계산 가속이 발생합니다.
        \end{enumerate}
    \end{qcblock}
    \vspace{0.2cm}
    \begin{qctip}{3일차 예고: Qiskit 회로 실습과 Grover 알고리즘}
        3일차부터는 글로벌 표준 프레임워크인 \textbf{Qiskit}을 도입하여 양자 회로를 직접 빌드하고,
        위상 간섭을 이용해 비정렬 검색 속도를 획기적으로 높이는 \textbf{Grover 알고리즘}을 탐구합니다.
    \end{qctip}
\end{frame}

\end{document}